%%%%%%%%%%%%%%%%%%%%%%%%%%%%%%%%%%%%%%%%%%%%%%%%%%%%%%%%%%%%%%%%%%%
% Frontiers Supplementary Material — AI prompts disclosure
% Class: frontiers_suppmat
% Per Frontiers BE WISE guidance (April 2026): list the generative-AI
% prompts used in manuscript preparation.
%%%%%%%%%%%%%%%%%%%%%%%%%%%%%%%%%%%%%%%%%%%%%%%%%%%%%%%%%%%%%%%%%%%

\documentclass[utf8]{frontiers_suppmat}
\usepackage{url,hyperref,lineno,microtype}
\hypersetup{hidelinks}
\usepackage[onehalfspacing]{setspace}

\begin{document}
\onecolumn
\firstpage{1}

\title[Supplementary Material]{{\helveticaitalic{Supplementary Material}}}

\maketitle

\section{Generative AI Prompts Used in Manuscript Preparation}

{\bf Manuscript:} Diagnostics and Infrastructure for Foundation Model-Based
Multi-Agent Systems: A Review of the Runtime Stack

{\bf Journal:} Frontiers in Computer Science

{\bf Date:} 2026-09-11

\subsection{Scope of this disclosure}

Per the Frontiers BE WISE AI guidance (April 2026), this file lists the
generative-AI prompts used in the preparation of the above manuscript.
AI-assisted tools were used for four tasks:

\begin{enumerate}
  \item \textbf{Literature search and screening} -- identifying candidate
        papers within the survey's scope of multi-agent diagnostics and
        infrastructure.
  \item \textbf{Taxonomy construction and organization} -- structuring the
        reviewed literature into the diagnostic and infrastructural
        categories reported in the manuscript.
  \item \textbf{Draft generation} -- producing initial prose for the
        manuscript sections, subsequently revised by the authors.
  \item \textbf{Reference verification} -- cross-checking bibliographic
        metadata (titles, authors, venues, publication years, arXiv
        identifiers) against the source papers.
\end{enumerate}

The prompts below are the canonical/representative set used for each task.
Long-running interactive refinement sessions produced many variations and
follow-up prompts; the categories and purposes listed here are the
substantive ones. No prompts were used to fabricate results, data, or
quantitative claims. Quantitative claims used for cross-system comparison
were verified against the primary sources; results marked [a] in the
manuscript are reported as stated in source abstracts and are not
independently verified.

The authors take full responsibility for the content, accuracy, and
completeness of the manuscript.

\subsection{Literature Search and Screening}

Purpose: identify candidate papers on multi-agent system diagnostics and
infrastructure, and filter to those relevant to the survey's scope.

Representative prompts:
\begin{itemize}
  \item ``Search for recent papers (2024--2026) on failure localization in
        multi-agent LLM systems. For each, report title, authors, venue,
        year, and arXiv ID if available.''
  \item ``Find papers on observability and tracing standards for LLM
        agents, including OpenTelemetry generative-AI semantic conventions
        and agent trace provenance frameworks.''
  \item ``Search for infrastructure and serving systems that optimize GPU
        memory, KV-cache management, or scheduling specifically for
        multi-agent workloads, not single-model inference.''
  \item ``Identify safety guardrail systems for multi-agent LLM systems:
        prompt injection defenses, policy verification of agent traces,
        and topology-aware defenses.''
  \item ``List benchmarks for evaluating agent reliability, failure
        attribution, and safety in multi-agent systems (e.g., trace-level,
        step-level, and safety-specific benchmarks).''
\end{itemize}

\subsection{Taxonomy Construction and Organization}

Purpose: organize the identified literature into coherent categories and
subcategories for the survey.

Representative prompts:
\begin{itemize}
  \item ``Group these failure localization methods into paradigms
        (statistical, LLM-guided, RL-trained). For each paradigm, note the
        input requirements (number of executions) and attribution
        granularity (agent/step/action).''
  \item ``Organize observability and tracing work into three levels:
        conceptual provenance models, standardization protocols, and
        production tooling.''
  \item ``Categorize multi-agent safety systems by abstraction level:
        topology-based, policy verification, and ML-layered defenses.''
  \item ``Structure multi-agent infrastructure systems by optimization
        axis: memory management, prefix-aware scheduling, cache lifecycle,
        and power management. Identify which workload characteristic each
        axis addresses.''
  \item ``Identify gaps: which dimensions are under-covered by existing
        benchmarks, and where diagnostics and infrastructure systems could
        be co-designed.''
\end{itemize}

\subsection{Draft Generation}

Purpose: produce initial prose for manuscript sections, subsequently edited
and verified by the authors.

Representative prompts:
\begin{itemize}
  \item ``Write an abstract for a survey covering two contributions: a
        diagnostics taxonomy (failure localization, observability, safety,
        benchmarks) and an infrastructure challenge analysis (memory,
        cache lifecycle, scheduling, power).''
  \item ``Draft an introduction establishing the gap: the diagnostics and
        infrastructure layers of multi-agent systems lack systematic joint
        analysis. Include motivation from adoption data.''
  \item ``Write the failure localization subsection, synthesizing
        statistical, LLM-guided, and RL-trained approaches as three
        paradigms, citing the provided papers.''
  \item ``Write the infrastructure section, presenting multi-agent serving
        challenges (sparse activation, invocation distance, shared
        prefixes) and application-aware systems that address each.''
  \item ``Write the discussion: open challenges, diagnostics-infrastructure
        co-design opportunities, broader impact, limitations, and future
        work.''
\end{itemize}

\subsection{Reference Verification}

Purpose: check bibliographic metadata against authoritative sources to
avoid citation errors and ensure all cited works are real and correctly
attributed.

Representative prompts:
\begin{itemize}
  \item ``For each of the following arXiv IDs, return the exact title,
        author list, venue, and publication year as listed in the source
        metadata. Flag any that do not resolve.''
  \item ``Verify the venue and publication year for these conference papers
        (ICLR, ICML, ACL, NeurIPS, COLM, EACL). Correct any year or venue
        mismatches.''
  \item ``Check these references for withdrawn status or known retractions.
        Note any paper whose authors have withdrawn or retracted it.''
  \item ``Confirm the numeric claims attributed to these papers appear in
        the source (e.g., reported accuracy, speedup, reduction in attack
        success rate) and report the exact value as published.''
\end{itemize}

\end{document}
